\subsection{Распространение излучения в ОВ}
Распространение излучение в ОВ, как и в прочих средах, описывается уравнениями Максвелла~\cite{kalluri2013principles}. В СИ данные уравнения имеют вид:

\begin{subequations}\label{eq:maxwell}
	\begin{align}
		\nabla \times \mathbf{E} &= -\frac{\partial \mathbf{B}}{\partial t}, \\
		\nabla \times \mathbf{H} &= \mathbf{J} + \frac{\partial \mathbf{D}}{\partial t}, \\
		\nabla \cdot \mathbf{D} &= \rho_f, \\
		\nabla \cdot \mathbf{B} &= 0.
	\end{align}
\end{subequations}

В системе~\cref{eq:maxwell} величины $\mathbf{E}$ и $\mathbf{H}$~--- вектора напряжённости электрического и магнитного полей соответственно,  $\mathbf{D}$ и $\mathbf{B}$~--- соответствующие векторы электрической и магнитной индукции. Вектор плотности тока $\mathbf{J}$ и плотность свободных зарядов $\rho_f$ описывают источники электромагнитного поля. Если в среде отсутствуют свободные заряды, эти величины обращаются в нуль: $\mathbf{J}=\rho_f=0$.

Вектора индукции $\mathbf{D}$ и $\mathbf{B}$ возникают как отклик среды на распространение электрического и магнитного полей E и H. Связь между ними задаётся материальными уравнениями:

\begin{subequations}\label{eq:material}
	\begin{align}
		\mathbf{D} &= \epsilon_0 \mathbf{E} + \mathbf{P}, \\
		\mathbf{B} &= \mu_0 (\mathbf{H} + \mathbf{M}).
	\end{align}
\end{subequations}

В системе материальных уравнений~\cref{eq:material} $\epsilon_0$~--- электрическая постоянная, $\mu_0$~--- магнитная постоянная, $\mathbf{P}$ и $\mathbf{M}$~--- вектора индуцированной электрической и магнитной поляризации соответственно. Для немагнитной среды, такой как ОВ, $\mathbf{M} = 0$.

Исключая переменные $\mathbf{B}$ и $\mathbf{D}$ из \cref{eq:maxwell, eq:material}, находим связь между $\mathbf{E}$ и $\mathbf{H}$:

\begin{equation}
	\nabla \times \nabla \times \mathbf{E} = -\frac{1}{c^2} \frac{\partial^2 \mathbf{E}}{\partial t^2} - \mu_0 \frac{\partial^2 \mathbf{P}}{\partial t^2} \label{eq:wave}
\end{equation}

Здесь $c$~--- скорость света в вакууме, $\mu_0\epsilon_0 = 1/c^2$. Для решения волнового уравнения~\cref{eq:wave} необходима связь между поляризацией $\mathbf{P}$ и напряжённостью электрического поля $\mathbf{E}$. Ограничиваясь нелинейностями третьего порядка ($\chi^{(3)}$), индуцированную поляризацию можно представить в виде суммы линейного и нелинейного (по электрическому полю) слагаемых~\cite{boyd2008nonlinear}:

\begin{equation}
	\mathbf{P} = \varepsilon_0 \left( \chi^{(1)} \cdot \mathbf{E} + \chi^{(2)} : \mathbf{EE} + \chi^{(3)} \vdots \mathbf{EEE} + \ldots \right) = P_{\mathrm{L}} + P_{\mathrm{NL}}
	\label{eq:polarization}
\end{equation}

В~\cref{eq:polarization} предполагается мгновенный отклик среды, поэтому данное разложение не описывает эффекты, связанные с запаздыванием, такие как ВКР и ВРМБ. Восприимчивость первого порядка $\chi^{(1)}$ соответствует линейному отклику среды и определяет комплексный показатель преломления среды. Восприимчивость первого порядка $\chi^{(2)}$ отвечает за такие н-о эффекты, как ГВГ, ГСЧ, ГРЧ, оптическое выпрямление. В средах с отсутствием центра инверсии (таких как кварцевое стекло в ОВ) данная величина равна нулю. Однако, как будет показано в Главе~\cref{ch:nonlin2}, в ОВ возможна фотоиндуцированная запись решётки квадратичной нелинейности.